\documentclass{resume}
\usepackage[left=0.75in,top=0.6in,right=0.75in,bottom=0.6in]{geometry}

% Chinese support (Overleaf-safe)
\usepackage[UTF8,fontset=fandol]{ctex}

\usepackage{bibentry}
\usepackage{url}
\newcounter{qcounter}

%----------------------------------------------------------------------------------------
% 联系方式
%----------------------------------------------------------------------------------------
\name{郦乐宁}
\address{美国 马萨诸塞州 波士顿}
\address{+1 (774) 823 2639 \\ leningli@outlook.com}

\begin{document}
\small

\nobibliography{refs}
\bibliographystyle{ieeetr}

%========================================================================================
% 研究兴趣
%========================================================================================
\begin{rSection}{研究兴趣}
面向机器人的智能体人工智能；强化学习；博弈论；形式化方法
\end{rSection}

%========================================================================================
% 教育背景
%========================================================================================
\begin{rSection}{教育背景}

{\bf 卡内基梅隆大学（Carnegie Mellon University, CMU）} \hfill 美国 宾夕法尼亚州 匹兹堡 \\
{\bf 工商管理硕士（MBA）} \hfill 2025.08 -- 至今 {\em（预计 2028.05 毕业）}

{\bf 伍斯特理工学院（Worcester Polytechnic Institute, WPI）} \hfill 美国 马萨诸塞州 伍斯特 \\
{\bf 机器人工程博士（Ph.D.）} \hfill 2016.08 -- 2022.12
\begin{list}{$\circ$}{\leftmargin=1.5em}
\item 博士论文：\textit{基于时间逻辑约束的随机系统最优控制与强化学习}
\end{list}

{\bf 机器人工程硕士（M.S.）} \hfill 2016.08 -- 2018.05 \\

{\bf 计算机科学硕士（M.S.）} \hfill 2014.08 -- 2016.05
\begin{list}{$\circ$}{\leftmargin=1.5em}
\item 硕士论文：\textit{BiRRTOpt：面向 Atlas 机器人的运动规划统一软件框架}
\end{list}

{\bf 哈尔滨工业大学（Harbin Institute of Technology, HIT）} \hfill 中国 \\
{\bf 信息安全学士（B.E.）} \hfill 2010.09 -- 2014.07
\begin{list}{$\circ$}{\leftmargin=1.5em}
\item 荣誉：Summa Cum Laude（专业前 5\%）
\item 本科论文：\textit{基于 Contourlet 变换的图像压缩}
\end{list}

{\bf 英语语言与文学学士（B.A.）} \hfill 2011.09 -- 2014.07
\begin{list}{$\circ$}{\leftmargin=1.5em}
\item 本科论文：\textit{《恋爱中的女人》中的男性沙文主义研究}
\end{list}

\end{rSection}

%========================================================================================
% 行业与学术顾问经历
%========================================================================================
\begin{rSection}{行业与学术顾问经历}

\begin{rSubsection}{研究附属人员（Research Affiliate）}{2026.10 -- 至今}{耶鲁大学}{美国 康涅狄格州 纽黑文}
\item 与 Na Li 教授合作，研究面向机器人的智能体人工智能（Agentic AI）。
\end{rSubsection}

\begin{rSubsection}{机器人实验室顾问}{2025.12 -- 至今}{哈佛大学}{美国 马萨诸塞州 奥尔斯顿}
\item 指导学生开展算法设计、系统评估与可复现的机器人实验。
\item 指导强化学习、运动规划、控制与安全关键系统方面的研究项目。
\item 评审研究计划、明确项目范围，并提供技术反馈。
\end{rSubsection}

\begin{rSubsection}{高级软件工程师}{2022.10 -- 2026.09}{Symbotic}{美国 马萨诸塞州 威明顿}
\item 使用 C++ 设计并部署面向数千台仓储机器人的路径规划与协同算法。
\item 通过改进控制设计、状态估计与容错行为架构，提升机器人集群的鲁棒性。
\item 担任远程操作系统技术负责人，协调研发与运营团队，将系统部署到生产环境。
\end{rSubsection}

\begin{rSubsection}{高级软件工程师}{2021.10 -- 2022.08}{Berkshire Grey}{美国 马萨诸塞州 贝德福德}
\item 负责未知 SKU 项目，使用 ROS、C++ 和 Python，使机器人能够抓取此前未见过的商品。
\item 通过重新设计进程间通信与执行流水线，降低端到端系统延迟。
\end{rSubsection}

\begin{rSubsection}{软件工程师实习}{2015.06 -- 2016.01}{Rudolph Technologies}{美国 马萨诸塞州 特克斯伯里}
\item 开发工具，将多套遗留代码库自动迁移至统一平台。
\item 改进采集与分析算法，提高晶圆缺陷数据的准确性。
\end{rSubsection}

\begin{rSubsection}{软件工程师实习}{2013.07 -- 2013.08}{东软集团（Neusoft）}{中国}
\item 开发地图管理系统，支持地理数据的插入、删除与更新。
\end{rSubsection}

\end{rSection}

%========================================================================================
% 代表性技能
%========================================================================================
\Needspace{12\baselineskip}
\begin{rSection}{代表性技能}
\begin{tabular}{ @{} >{\bfseries}l @{\hspace{3ex}} l }
人工智能与感知： & 强化学习，PyTorch，TensorFlow，OpenCV \\
机器人与系统： & ROS / ROS 2，Linux，多机器人系统，远程操作 \\
编程语言： & Python，C++，C，MATLAB \\
语言能力： & 英语（流利），中文（母语）
\end{tabular}
\end{rSection}

%========================================================================================
% 教学经历
%========================================================================================
\begin{rSection}{教学经历}

\begin{rSubsection}{助教}{2022.08 -- 2022.12}{RBE 549 计算机视觉}{伍斯特理工学院}
\item 共同设计并讲授新开研究生计算机视觉课程，涵盖经典方法与深度学习。
\item 负责每周讲课、答疑与学生项目指导。
\item 课程网站：https://nitinjsanket.github.io/teaching/rbe549/fall2022.html
\end{rSubsection}

\begin{rSubsection}{助教}{2020.08 -- 2021.05}{RBE 3001 \& 3002 统一机器人系统 III \& IV}{伍斯特理工学院}
\item 指导 3D 打印机械臂运动学/控制与自主移动机器人导航实验。
\item 设计并批改课程项目、实验报告与作业。
\end{rSubsection}

\begin{rSubsection}{助教}{2018.08 -- 2018.12}{RBE 549 计算机视觉}{伍斯特理工学院}
\item 为研究生讲授计算机视觉理论与实现，并主持课程答疑。
\end{rSubsection}

\begin{rSubsection}{助教}{2017.08 -- 2018.05}{RBE 1001 机器人导论}{伍斯特理工学院}
\item 管理并指导 5 名本科生助教。
\item 设计并批改课程项目、实验报告与作业。
\end{rSubsection}

\end{rSection}

%========================================================================================
% 代表性研究项目
%========================================================================================
\begin{rSection}{代表性研究项目}

\begin{rSubsection}{Kinex：面向编程智能体的机器人框架}{2026 -- 至今}{}{\url{https://stacx-lab.github.io/kinex/}}
\item 参与开发开源机器人框架，将编程智能体连接到真实机器人与仿真环境，并支持 ROS 2 集成。
\item Kinex 的技能、工具与设备架构支持代码执行与机器人协同，并通过观测和执行反馈指导智能体行动。
\item 1.0.0 版本于 2026 年 10 月 15 日发布。源代码：\url{https://github.com/stacx-lab/kinex}。
\end{rSubsection}

\begin{rSubsection}{DARPA Robotics Challenge}{2014.09 -- 2015.06}{}{美国 马萨诸塞州 伍斯特}
\item 与 WPI-CMU 团队合作，在 DARPA Robotics Challenge 决赛中参与波士顿动力（Boston Dynamics）Atlas 人形机器人项目。
\item 设计用于开门、阀门操作与工具抓取的机械臂运动规划算法。
\item 设计操作员界面，支持态势感知与任务执行。
\item 所在 WPI-CMU 团队在决赛中获 24 支参赛队伍中的第 7 名。
\end{rSubsection}

\begin{rSubsection}
{DARPA SI3-CMD：复杂军事决策中的不完美信息序列交互}
{2019.01 -- 2020.08}{}{美国 马萨诸塞州 伍斯特}
\item 与 SSCI 合作，开发基于博弈论的欺骗性规划框架及 Python 实现。
\item 在形式化超博弈模型中利用信息不对称与策略性欺骗，提升任务成功概率。
\item 提出面向时间逻辑目标的动态超博弈（hypergame）解概念。
\end{rSubsection}

\begin{rSubsection}
{基于时间逻辑约束的随机系统最优控制与强化学习}
{2016.08 -- 2022.08}{}{美国 马萨诸塞州 伍斯特}
\item 构建将概率时间逻辑规范转化为具满足性保证的机会约束控制问题的理论框架。
\item 提出适用于连续随机系统的可扩展无模型强化学习方法，提高样本效率。
\end{rSubsection}

\end{rSection}

%========================================================================================
% 荣誉与奖励
%========================================================================================
\begin{rSection}{荣誉与奖励}

\begin{nSubsection}
{Tepper 职业发展基金}{2026}
{卡内基梅隆大学}{美国 宾夕法尼亚州 匹兹堡}
\\ 4 月 9 日及 9 月 30 日批准。
\end{nSubsection}

\begin{nSubsection}
{Tepper 商学院在线混合 MBA 奖学金}{2025}
{卡内基梅隆大学}{美国 宾夕法尼亚州 匹兹堡}
\end{nSubsection}

\begin{nSubsection}
{Alex F. Backlin 基金奖学金}{2021.01}
{伍斯特理工学院}{美国 马萨诸塞州 伍斯特}
\end{nSubsection}

\begin{nSubsection}
{差旅资助奖}{2019.10}
{利哈伊大学}{美国 宾夕法尼亚州 伯利恒}
\end{nSubsection}

\begin{nSubsection}
{研究生差旅奖}{2019.10}
{伍斯特理工学院}{美国 马萨诸塞州 伍斯特}
\end{nSubsection}

\begin{nSubsection}
{研究生差旅奖}{2019.03}
{伍斯特理工学院}{美国 马萨诸塞州 伍斯特}
\end{nSubsection}

\begin{nSubsection}
{研究生差旅奖}{2017.06}
{伍斯特理工学院}{美国 马萨诸塞州 伍斯特}
\end{nSubsection}

\end{rSection}

%========================================================================================
% 资格认证
%========================================================================================
\begin{rSection}{资格认证}

\begin{hSubsection}
{REC 基金会认证教练}
{2024.09 -- 2025.09}
{Robotics Education \& Competition Foundation}
{线上}
{2024--2025 赛季执教 Vinci STEAM 初中 VEX 队 10102Z；该队获 2025 年 VEX 机器人世界锦标赛（达拉斯）初中组世界冠军。}
\end{hSubsection}

\begin{hSubsection}
{FIRST Tech Challenge（FTC）认证教练}
{2023.11 -- 2024.05}
{FIRST Robotics}
{线上}
{指导 FTC 机器人团队进行机器人设计、编程与竞赛策略制定。}
\end{hSubsection}

\begin{hSubsection}
{大学教学资格认证}
{2017.06 -- 2019.08}
{马萨诸塞州中部高等教育联盟（HECCMA）}
{美国 马萨诸塞州 伍斯特}
{完成循证教学法与高校教学培训。}
\end{hSubsection}

\end{rSection}

%========================================================================================
% 课外活动与领导力
%========================================================================================
\Needspace{10\baselineskip}
\begin{rSection}{课外活动与领导力}

\begin{rSubsection}{主席}{2019.01 -- 2020.05}{WPI 研究生学生会（Graduate Student Government, GSG）}{美国 马萨诸塞州 伍斯特}
\item 领导 WPI 研究生学生会，并协调学生会治理事务。
\item 担任研究生群体与学校管理层之间的主要沟通桥梁。
\item 与研究生院合作，向董事会反映研究生住房问题。
\item 受邀加入教务长（Provost）遴选委员会，并提交正式建议。
\end{rSubsection}

\begin{rSubsection}{志愿者}{2013.07 -- 2013.09}{拉萨儿童福利院}{中国}
\item 为低收入家庭儿童的教育服务项目筹集资金。
\item 为儿童提供语文、数学与英语辅导。
\end{rSubsection}

\end{rSection}

%========================================================================================
% 学术与专业组织成员
%========================================================================================
\begin{lSection}{学术与专业组织成员}
\item {\bf 会员}，IEEE
\item {\bf 会员}，IEEE Young Professionals
\item {\bf 会员}，IEEE Robotics and Automation Society
\item {\bf 会员}，Rho Beta Epsilon 荣誉学会 Alpha 分会（伍斯特理工学院）
\end{lSection}

%========================================================================================
% 学术服务
%========================================================================================
\begin{rSection}{学术服务}

\begin{lSubsection}{期刊审稿人}
\item IEEE Transactions on Automation Science and Engineering (T-ASE)
\item IEEE Transactions on Systems, Man, and Cybernetics: Systems
\item IEEE Transactions on Intelligent Transportation Systems
\item IEEE Robotics and Automation Letters (RA-L)
\item IEEE Access
\item IET Cyber-Systems and Robotics
\item Discover Robotics
\item The Journal of Supercomputing
\end{lSubsection}

\begin{lSubsection}{会议审稿人}
\item International Conference on Robotics and Automation (ICRA)
\item American Control Conference (ACC)
\item IEEE Conference on Decision and Control (CDC)
\item IEEE/RSJ International Conference on Intelligent Robots and Systems (IROS)
\item International Conference on Ubiquitous Robots (UR)
\item European Control Conference (ECC)
\end{lSubsection}

\end{rSection}

%========================================================================================
% 学术主页与社交媒体
%========================================================================================
\begin{lSection}{学术主页与社交媒体}
\begin{tabular}{ @{} >{\bfseries}l @{\hspace{6ex}} l } \\
LinkedIn: & \url{https://www.linkedin.com/in/lening-li/} \\
个人主页: & \url{https://lening.li} \\
GitHub: & \url{https://github.com/leelening} \\
Google Scholar: & \url{https://scholar.google.com/citations?user=KWUJ10wAAAAJ\&hl=en}
\end{tabular}
\end{lSection}

%========================================================================================
% 论文发表（保持英文）
%========================================================================================
\begin{rSection}{论文发表}

\begin{pSubsection}{arXiv Preprints}{P.}
\item \bibentry{li2021policy}
\end{pSubsection}

\Needspace{6\baselineskip}
\begin{pSubsection}{Conferences}{C.}
\item \bibentry{li2023topological}
\item \bibentry{li2023probabilistic} {(录用率：15\%)}
\item \bibentry{li2023synthesis}
\item \bibentry{li2019topological}
\item \bibentry{li2019approximate}
\item \bibentry{li2017sampling}
\item \bibentry{li2016birrtopt}
\item \bibentry{atkeson2015no}
\end{pSubsection}

\begin{pSubsection}{Journals}{J.}
\item \bibentry{li2022dynamic}
\item \bibentry{chen2018building}
\item \bibentry{dedonato2017team}
\end{pSubsection}

\begin{pSubsection}{Chapters}{Ch.}
\item \bibentry{atkeson2018achieving}
\item \bibentry{atkeson2018happened}
\end{pSubsection}

\end{rSection}

\end{document}
